\documentclass[11pt,letterpaper]{article}
\PassOptionsToPackage{hyphens}{url}
\usepackage[margin=.83in,top=.8in,bottom=.8in,headheight=15pt]{geometry}
\usepackage[T1]{fontenc}
\usepackage{lmodern,amsmath,amssymb,mathtools,bm,booktabs,array,graphicx,tikz}
\usetikzlibrary{arrows.meta,positioning,calc,decorations.pathreplacing,patterns}
\definecolor{ink}{HTML}{17344D}\definecolor{accent}{HTML}{087E8B}\definecolor{pale}{HTML}{EDF5F5}\definecolor{gold}{HTML}{AD6B19}
\providecommand{\dd}{\,\mathrm{d}}\providecommand{\R}{\mathbb{R}}
\providecommand{\E}{\mathbb{E}}\providecommand{\Var}{\operatorname{Var}}
\providecommand{\Cov}{\operatorname{Cov}}\providecommand{\diag}{\operatorname{diag}}
\providecommand{\norm}[1]{\left\lVert#1\right\rVert}
\providecommand{\abs}[1]{\left|#1\right|}
\usepackage{microtype,hyperref,fancyhdr,enumitem,amsthm,sectsty,needspace}
\allsectionsfont{\raggedright}
\hypersetup{colorlinks=true,linkcolor=ink,urlcolor=accent,pdfauthor={Zan Ahmad},pdftitle={ODE integration, root finding, and parameter estimation / Review}}
\setlength{\parindent}{0pt}\setlength{\parskip}{6pt}\linespread{1.07}
\setlist{itemsep=3pt,topsep=4pt,leftmargin=*}
\emergencystretch=2em\clubpenalty=10000\widowpenalty=10000
\pagestyle{fancy}\fancyhf{}\fancyhead[L]{\small Scientific Computing in Python}\fancyhead[R]{\small Lecture 4}
\fancyfoot[L]{\footnotesize Zan Ahmad}\fancyfoot[R]{\footnotesize\thepage}\renewcommand{\headrulewidth}{.3pt}
\newtheorem{proposition}{Proposition}\newtheorem{theorem}{Theorem}\newtheorem{lemma}{Lemma}
\theoremstyle{definition}\newtheorem{definition}{Definition}\newtheorem{example}{Example}
\numberwithin{equation}{section}
\begin{document}
\section*{ODE integration, root finding, and parameter estimation --- review sheet}
Use this after attempting the worked examples and practice. Explain each claim in words before calculating.

\section*{Key ideas}
\begin{itemize}
\item One model can support several distinct numerical questions.
\item A rate law and initial value define an initial-value problem.
\item Exact-model comparisons check numerical implementation, not real-world validity.
\item A bracketed root needs a suitable continuous residual and endpoint evidence.
\item Interpolation adds a between-sample assumption and has a declared domain.
\item Optional optimization fits a parameter within bounds; residuals and independent evidence still matter.
\end{itemize}
\section*{Common errors and repairs}
\textbf{Giving k units of minutes.} k must have reciprocal-minute units so k times t is dimensionless.

\textbf{Treating solver tolerance as a guaranteed global error bound.} Measure the relevant error against an independent reference when available.

\textbf{Calling a target equal to ambient a finite crossing.} The exponential reaches that limit only asymptotically.

\textbf{Treating more interpolation plot points as new measurements.} The same knots and linear assumption define all those plotted values.

\textbf{Treating the best tested parameter as model validation.} A smaller residual does not establish the physical model; compare independent observations.

\section*{Retrieve without your notes}
\begin{itemize}
\item What happens when k is zero?
\item Why compare a root with both an analytic time and a residual?
\item What is the interpolation domain?
\item What must be added to assess a model beyond synthetic recovery?
\end{itemize}
\section*{Readiness check}
\begin{itemize}
\item State a one-state cooling model with consistent time and temperature units.
\item Run a supplied initial-value solver and compare its output with an exact solution of the same model.
\item Translate a target-temperature question into a bracketed scalar root.
\item Interpolate inside a sampled interval and explain why interpolation is not the ODE model.
\item Use limiting behavior and residuals to check numerical outputs.
\end{itemize}
\end{document}
